\begin{abstract}
\noindent
Does US aid buy rhetorical allegiance? We test the Kuziemko--Werker (2006) hypothesis that UN Security Council membership, by attracting additional US foreign assistance, induces countries to express more pro-US sentiment in their UN General Debate speeches. Using an LLM-based semantic stance measure applied to 7,866 speeches from 179 countries (1970--2020), we find no evidence for this causal chain. The first stage fails under every fixed-effect specification: the $F$-statistic on UNSC membership in a country--year FE regression of log aid ranges from 0.1 to 1.1 and never approaches the Staiger--Stock $F>10$ threshold. The reduced form is consistently negative across specifications (country--year FE: $\hat{\beta} = -0.100$, $p=0.075$), trending toward \emph{less} pro-US speech among UNSC members. An event study reveals no systematic pre-trends or post-term effects. Even replicating Kuziemko--Werker's own region--year FE specification, the first stage is not statistically significant ($F=0.9$). We conclude that in this data, UNSC membership is not a valid instrument for US aid, and the ``material power buys allegiance'' thesis finds no empirical support through this channel. We discuss the fragility of the UNSC--aid IV chain and the gap between diplomatic rhetoric and observable behavior.

\vspace{1em}
\noindent\textbf{JEL Classification:} D72, F35, F51, O19, C26.

\noindent\textbf{Keywords:} UN Security Council, foreign aid, allegiance, UN General Debate, instrumental variables, text-as-data, LLM measurement.
\end{abstract}
